\documentclass[11pt]{article}
\usepackage[margin=1in]{geometry}
\usepackage{amsmath,amssymb}
\usepackage[hidelinks]{hyperref}

\newcommand{\lk}{\operatorname{lk}}
\newcommand{\Ftwo}{\mathbb{F}_2}

\title{A six-vertex counterexample to Conjecture 00000002729}
\author{}
\date{}

\begin{document}
\maketitle

The conjecture asserts that a simplicial complex is Cohen--Macaulay if and
only if it is connected and every vertex link is Cohen--Macaulay. A pure
two-dimensional complex disproves the converse over every coefficient field.

\section*{The complex and its links}

Let $\Delta$ consist of all subsets of the following six facets, where $abc$
denotes the set $\{a,b,c\}$:
\[
  014,\quad 034,\quad 125,\quad 145,\quad 023,\quad 235.
\]
Every facet is a triangle and all six vertices occur. The edges along
$0-1-2-3-4-5$ form a spanning path, so $\Delta$ is connected.
For each vertex, deleting that vertex from each facet containing it gives
the edges of its link. The resulting link graphs are exactly these paths:
\[
\begin{array}{c|c@{\qquad}c|c}
v & \lk_\Delta(v) & v & \lk_\Delta(v) \\ \hline
0 & 1-4-3-2 & 3 & 4-0-2-5\\
1 & 0-4-5-2 & 4 & 3-0-1-5\\
2 & 0-3-5-1 & 5 & 4-1-2-3
\end{array}
\]

By Reisner's criterion~\cite[Theorem~5]{reisner}, a finite simplicial complex
is Cohen--Macaulay over a field $k$ exactly when, for every face $F$ including
the empty face, $\widetilde H_i(\lk_\Delta(F);k)=0$ for
$i<\dim\lk_\Delta(F)$. In a pure one-dimensional complex, the empty-face
condition is connectedness, since the complex is nonempty. The link of each
vertex is zero-dimensional and has a vertex because each original vertex is
in an edge. The edge links have dimension $-1$ and impose no condition.
Consequently every path in the table is Cohen--Macaulay over $k$.

\section*{The obstruction for the empty face}

Write $[a,b]$ for an oriented edge. Define a $k$-valued linear functional $c$
on the oriented edge chains by its values on edges with increasing endpoints:
\[
 c([0,2])=1,\qquad c([2,3])=-1,\qquad c([3,5])=1,
\]
with zero values on all other increasing edges. The facets $014$, $034$,
$125$, and $145$ have no edges with nonzero value. The other two boundaries
evaluate as
\[
\begin{aligned}
c(\partial[0,2,3])&=c([2,3]-[0,3]+[0,2])=-1+1=0,\\
c(\partial[2,3,5])&=c([3,5]-[2,5]+[2,3])=1-1=0.
\end{aligned}
\]
Thus $c$ vanishes on every boundary of a two-chain. The one-chain
$z=[0,1]+[1,2]-[0,2]$ is a cycle, since each vertex has total boundary
coefficient zero, and $c(z)=-1\ne0$. Therefore $z$ is not a boundary and
$\widetilde H_1(\Delta;k)=H_1(\Delta;k)\ne0$ over every field $k$.
The empty face has link $\Delta$ of dimension two, so Reisner's criterion
rules out Cohen--Macaulayness.

\section*{Lean verification}

The standalone Lean 4 project checks the disproof over $\Ftwo$, with the same
triangles and the same cycle. Its predicate \texttt{cohenMacaulayF2} spells out
Reisner's conditions for a pure two-dimensional complex. For the empty face
these require connectedness and vanishing first homology. For a vertex, they
require a nonempty, connected one-dimensional link. For an edge, they require
its zero-dimensional link to have a vertex. A triangle link has no requirement
below dimension $-1$.

A list of edges or triangles represents a chain by counting the multiplicity
of each simplex modulo two. Every simplicial chain over $\Ftwo$ has such a
representation. The code verifies that the displayed cycle has zero boundary
and compares edge chains coefficient by coefficient. It then proves that
the sum of the three indicated edge coefficients modulo two vanishes on the
boundary of every two-chain, by induction on the list of triangles. This sum
is one on the cycle. The theorem \texttt{C2729.not\_localCriterion} negates
the conjectured equivalence already on complexes represented by six
triangular facets on six vertices. Run \texttt{lake build} from the
\texttt{lean} directory to check the theorem.

\begin{thebibliography}{1}
\bibitem{reisner}
G. A. Reisner, \emph{Cohen--Macaulay quotients of polynomial rings},
Advances in Mathematics \textbf{21} (1976), 30--49.
\href{https://doi.org/10.1016/0001-8708(76)90114-6}{DOI 10.1016/0001-8708(76)90114-6}.
\end{thebibliography}

\end{document}
